\section{Prüfer und Guard-Propagation}\label{sec:guards}

Offene Bedingungen werden in QAM auf zwei Ebenen geführt, die verschieden tief geprüft sind. Ebene~A ist die Statusvererbung in der Akte;
ihre Eigenschaften sind für jede akzeptierte Akte in Lean bewiesen. Ebene~B sind die Guards der Ausführungsmaschine; ihre abstrakte
Übergangssemantik ist in Lean bewiesen, die tatsächlich laufende Implementierung nicht. Wir halten beide Ebenen getrennt, weil eine
Vermischung genau die Art von Beförderung wäre, die QAM verhindern soll.

\subsection{Drei Prüfer für die Akte}\label{sec:pruefer}

Die Regeln (R1)--(R4) werden dreimal unabhängig angewandt:
\begin{enumerate}[nosep]
\item \texttt{qam.sty} beim Kompilieren (Abschnitt~\ref{sec:latex});
\item \texttt{qamcheck}, ein in Lean~4 geschriebenes Programm, auf der exportierten Akte. Es prüft jede Zeile gegen die vor ihr stehenden
und kontrolliert, dass \texttt{sources} genau die eingetragenen Hashes der Prämissen sind und kein Hash doppelt vorkommt:
\begin{lstlisting}
cd qam/lean && lake build && lake exe qamcheck ../examples/example.qam
\end{lstlisting}
\item \texttt{qamhash.py verify}, das alle Eintragshashes, den Aktenhash und den MD5-Wert der Quelle neu berechnet.
\end{enumerate}
Zur Gegenprobe gehören absichtlich fehlerhafte Akten, die abgewiesen werden müssen: eine Beförderung, eine Frame-Vermischung, eine falsche
Unabhängigkeitserklärung, eine falsche Abbildungsquelle, ein gebrochener Hash-Verweis, ein nachträglich veränderter Eintrag und eine von Hand
gelöschte Eingabe. Die letzte kann \TeX{} gar nicht erzeugen; sie stellt sicher, dass der Lean-Prüfer nicht nur nachvollzieht, was das
\TeX-Paket ohnehin verhindert.

\subsection{Ebene A: Statusvererbung, in Lean bewiesen}\label{sec:ebeneA}

Das Lean-Paket \texttt{qam/lean} benutzt nur den Lean-Kern, ohne Mathlib, und hängt von keiner Anwendung ab. Es enthält zwei Module.

\paragraph{Herleitungsbäume.} \texttt{QAM.Provenance} formalisiert Herleitungsbäume mit drei Knotenarten: Schritt (eine Prämisse),
Verbindung (zwei Prämissen) und Überführung über eine Abbildung. Ein Baum ist wohlgeformt relativ zu einer Menge deklarierter Abbildungen,
wenn kein Knoten stärker ist als seine Prämissen, Verbindungen nur Prämissen eines Frames verbinden und jede Überführung eine deklarierte
Abbildung mit passender Quelle benutzt.

\begin{satz}[\texttt{no\_promotion}, \texttt{proved\_only\_from\_proved}]
In einem wohlgeformten Baum hat die Konklusion höchstens den Status jeder Prämisse. Eine Konklusion mit Status \st{proved} ruht nur auf
bewiesenen Aussagen.
\end{satz}

\begin{satz}[\texttt{inputs\_tracked}, \texttt{inputs\_sourced}, \texttt{independent\_iff}]
Die Eingaben der Konklusion sind genau die Eingaben ihrer Prämissen. Eine Konklusion ist genau dann unabhängig von einer Eingabe, wenn
keine Prämisse sie verwendet.
\end{satz}

\begin{satz}[\texttt{frames\_declared}, \texttt{no\_maps\_same\_frame}]
In einem wohlgeformten Baum erreicht jede Prämisse den Frame der Konklusion über eine Kette deklarierter Abbildungen. Ohne deklarierte
Abbildungen liegen alle Prämissen im Frame der Konklusion.
\end{satz}

\paragraph{Akten.} \texttt{QAM.QamFile} ist der Prüfer für exportierte Akten. Für jede Zeilenart ist bewiesen, was eine erfolgreiche Prüfung
bedeutet (\texttt{derive\_sound}, \texttt{convert\_sound}, \texttt{indep\_sound}, \texttt{compare\_sound}, für die Hash-Verweise
\texttt{link\_derive\_sound}, \texttt{link\_convert\_sound}, \texttt{link\_leaf\_sound}). Entscheidend sind die Sätze über ganze Akten:

\begin{satz}[\texttt{checked\_good}, \texttt{checked\_no\_promotion\_record}, \texttt{checked\_proved\_record}]
Jede Akte, die der Prüfer annimmt, erfüllt eine Invariante über alle Einträge. Insbesondere hat kein Eintrag einen höheren Status als ein
Eintrag, der von ihm aus über Prämissen in beliebig vielen Schritten erreichbar ist, und ein als \st{proved} geführter Eintrag ruht in
beliebiger Tiefe nur auf bewiesenen Einträgen.
\end{satz}

\begin{satz}[\texttt{checked\_inputs\_record}, \texttt{checked\_independent\_record}, \texttt{checked\_inputs\_sourced}]
In einer angenommenen Akte geht keine Eingabe entlang einer Prämissenkette verloren, jede Eingabe eines Eintrags stammt aus einer Prämisse,
und ein als unabhängig erklärter Eintrag ruht in keiner Tiefe auf einem Eintrag mit Eingabe.
\end{satz}

Das ist die formale Form der Vererbung offener Bedingungen aus Abschnitt~\ref{sec:offen}: \emph{In einer angenommenen Akte kann eine
offene Bedingung nicht dadurch verschwinden, dass weitere Schlüsse auf ihr aufbauen.} Die Sätze sind elementar, und das ist beabsichtigt.
Ihr Wert liegt nicht in mathematischer Tiefe, sondern darin, dass die Regeln präzise genug sind, um von einem Beweisassistenten und vom
\TeX-Compiler gleichermaßen geprüft zu werden.

\paragraph{Was Ebene A nicht prüft.} Nur Deklariertes wird geprüft. Der Parser, der die Textzeilen einliest, ist nicht verifiziert; Lean
berechnet kein MD5, das übernimmt das Skript. Ein ohne Prämissen als \st{proved} deklarierter Eintrag wird angenommen.

\subsection{Ebene B: Guards der Ausführungsmaschine}\label{sec:ebeneB}

Rechnungen, Suchen und Simulationen laufen nicht als Fließtext, sondern als Aufträge an eine Maschine (Abschnitt~\ref{sec:maschine}). Für
diese Ebene ist der Guard ein Ausgang eines Schritts.

\paragraph{Abstrakte Semantik, in Lean bewiesen.} Ein \emph{Speicher} besteht aus Inhalt (Objekte an Adressen, Relationen zwischen
Adressen), einem Guard-Register (Prädikate, die der Inhalt erfüllen muss), Schnappschüssen und einem Protokoll. Ein Speicher ist
\emph{zulässig}, wenn jede Relation gespeicherte Adressen benutzt und jeder Guard gilt. Jede Anweisung hat eine Rohwirkung, die undefiniert
sein kann, und ein Schritt wird so ausgeführt:
\[
\mathrm{exec}(s,i)=\begin{cases}
(s\ \text{mit Eintrag }(i,\textsc{reject}),\ \textsc{reject}) & \text{Rohwirkung undefiniert},\\
(s'\ \text{mit Eintrag }(i,\textsc{cont}),\ \textsc{cont}) & \text{Rohwirkung } s' \text{ zulässig},\\
(s\ \text{mit Eintrag }(i,\textsc{guard}),\ \textsc{guard}) & \text{sonst.}
\end{cases}
\]
Ein Ergebnis, das einen Guard verletzt, wird nie übernommen; der alte Inhalt bleibt, das Protokoll hält den Befund fest. Für dieses Modell
ist in Lean bewiesen (Anhang~\ref{app:lean}): Zulässigkeit bleibt erhalten (\texttt{exec\_preserves}, \texttt{run\_preserves}); kein Schritt
bleibt stecken (\texttt{exec\_ne\_halt}); \textsc{cont} und \textsc{reject} sind exakt charakterisiert (\texttt{exec\_cont\_iff},
\texttt{exec\_reject\_iff}); ein Schritt ohne \textsc{cont} lässt Inhalt, Schnappschüsse und Guards unverändert (\texttt{exec\_atomic}); jeder
Schritt erzeugt genau einen Protokolleintrag (\texttt{exec\_logs}, \texttt{run\_ledger\_length}); und es gibt kein Zertifikat für eine falsche
Aussage, gleich welche Anweisungen liefen (\texttt{no\_false\_certificate}). Ein Kandidat wird nur durch einen von außen gelieferten Beweis
zum Zertifikat; keine Anweisung erzeugt eines.

\paragraph{Implementierung, nicht verifiziert.} Die laufende Maschine ist ein Python-Adapter über einer SQLite-Datenbank. Sie folgt
derselben Dreiteilung: Jeder Auftrag endet mit \texttt{continue}, \texttt{guard} oder \texttt{reject}. Vor der Ausführung prüft sie die
deklarierten Guards von Kontext, Auftrag und beteiligten Objekten sowie die verknüpften Prämissen der Akte. Eine fehlende Prämisse ergibt
\texttt{guard} und verhindert die Ausführung. Ein Beleg mit \texttt{guard} wird mit seinen offenen Guards und Verpflichtungen
gespeichert und ist später abrufbar; er ist ein vollwertiges Ergebnis, kein Fehlschlag. Dass der Adapter die abstrakte Semantik korrekt
umsetzt, ist nicht bewiesen, sondern durch Tests und für die Bandmaschine durch einen Vergleich mit zwei Hardwarebeschreibungen gestützt
(Abschnitt~\ref{sec:maschine}).

\paragraph{Weitergabe.} Guards werden auf dieser Ebene auf drei Wegen weitergegeben:
\begin{enumerate}[nosep]
\item \emph{Prämissenspur.} Ein Auftrag kann Hashes von Akteneinträgen als Prämissen nennen. Die Maschine folgt deren Quellhashes
rekursiv und legt eine Momentaufnahme aller erreichten Einträge mit Wortlaut, Status und Frame im Beleg ab. Der Status bleibt dabei
erhalten: Eine als \st{reading} verknüpfte Zählregel bleibt eine Lesart, auch wenn die Rechnung auf ihr erfolgreich läuft.
\item \emph{Abhängige Aufträge.} In einem Stapel kann ein Auftrag von anderen abhängen (\texttt{depends\_on}). Endet ein Vorgänger mit
\texttt{guard}, \texttt{reject} oder ist selbst gesperrt, wird der Nachfolger nicht ausgeführt, sondern als \texttt{blocked} mit
\texttt{blocked\_by} geführt. Mit \texttt{requires\_match} läuft ein Nachfolger nur, wenn sein Vorgänger mindestens einen Kandidaten gefunden
hat; ein Ergebnis ohne Kandidaten bleibt dabei eine vollständige Rechnung.
\item \emph{Rückwärtssuche.} Wird ein Akteneintrag korrigiert, findet die Abfrage nach betroffenen Belegen alle gespeicherten Belege, die
ihn direkt, indirekt oder als fehlende Prämisse verwendet haben. Alte Belege werden nicht verändert.
\end{enumerate}

\subsection{Warum isolieren statt glätten}

Ein Prüfsystem, das nur zwei Ausgänge kennt, muss jede unentschiedene Frage einem der beiden zuschlagen. Nach Turing und Rice
(Abschnitt~\ref{sec:einleitung}) gibt es solche Fragen grundsätzlich, und in der Forschung ist die Mehrzahl der interessanten Aussagen für
lange Zeit unentschieden. Der Guard ist der dritte Ausgang, und QAM verlangt, dass er so behandelt wird wie die beiden anderen: gespeichert,
gehasht, abfragbar und vererbt.

Für die Arbeit mit Sprachmodellen hat das eine bestimmte Folge. QAM behandelt eine Halluzination nicht als inhaltlichen Fehler, den es
erkennen müsste; dazu wäre es nicht in der Lage. Es behandelt sie als \emph{Statusfehler}: als Versuch, eine Aussage stärker zu führen, als ihre
Prämissen erlauben, oder eine Eingabe als Ergebnis auszugeben. Diese Fehlerart ist mechanisch erkennbar, weil sie sich in der Struktur
der Akte zeigt und nicht in der Bedeutung des Texts. Ein Sprachmodell kann eine Bedingung im Fließtext weglassen; es kann sie nicht aus der
Akte entfernen, ohne dass sich Hashes ändern, und es kann den Status eines Schlusses nicht über den seiner Prämissen setzen, ohne dass der
Compiler anhält.
